\section{Related Work}


\textbf{LLM inference.} LLM inference has recently attracted significant research attention, resulting in a variety of systems designed to improve the efficiency and scalability of LLM inference. For instance, Orca~\cite{yu2022orca} introduces continuous batching to enhance throughput, while vLLM~\cite{vllm2024} proposes paged-attention for more granular KV cache management. Furthermore, both online and offline inference scenarios have been extensively studied. For online serving, DistServe~\cite{zhong2024distserve} proposes PD disaggregation to eliminate interference between the prefill and decode stages, and Sarathi~\cite{agrawal2024taming} implements a chunked prefill strategy to reduce contention between these operations. For offline inference, BlendServe~\cite{zhao2024blendserveoptimizingofflineinference} and BatchLLM ~\cite{zheng2025batchllmoptimizinglargebatched} leverage request reordering to accelerate processing for requests with shared prefixes and maximize resource utilization. 
Notably, TD-Pipe focuses on pipeline optimization to reduce bubbles, which is orthogonal to the request reordering method. \\
\textbf{Bubble Elimination for Pipeline Parallelism.} Pipeline parallelism is commonly employed in deep learning training and inference because it offers substantially lower communication overhead compared to tensor parallelism. However, complex dependency issues can give rise to pipeline bubbles, which ultimately lead to suboptimal device utilization. In terms of training, frameworks such as GPipe~\cite{huang2019gpipe}, Megatron ~\cite{narayanan2021efficient}, and Mobius~\cite{feng2023mobius} have introduced techniques to mitigate pipeline bubbles during the forward and backward passes. For inference, Sarathi~\cite{agrawal2024taming} initially proposed chunked prefill specifically to reduce pipeline bubbles.


% BlendServer, BatchLLM
% Prefix Cache
